\chapter{Methodology}
\label{chap:methodology}

This chapter describes how HepatoTwin is built and how it is evaluated.
Section~\ref{sec:m_data} introduces the data: the NHANES FibroScan cohort,
the patients' 24-hour dietary recalls, the recipe corpus and the external
NHANES~III cohort. Sections~\ref{sec:m_layer2} to \ref{sec:m_summary} mirror the different levels in
Figure~\ref{fig:hepatotwin_overview} namely the model of liver state and its safety criteria (Layer~2), 
the evidence engine (Layer~3), the recipe evaluation and prioritization engine (Layer~4), the scenario engine (Layer~5), 
the loop for updating the twin based on incoming data, and the layperson's report. Section~\ref{sec:m_system} describes the software
system and the patient and clinician interfaces (Layers~6 and~7).
Section~\ref{sec:m_eval} sets out how each component is evaluated, including
the in-silico evaluation of the twin on virtual patients. The experiments
behind the liver-state model are reported separately in
Chapter~\ref{chap:experiments}; the results of every evaluation are given in
Chapter~ef{chap:results}.

% ======================================================================
\section{Data}
\label{sec:m_data}

\subsection{NHANES FibroScan cohort}

The twin is built on the US National Health and Nutrition Examination Survey
(NHANES), which performed vibration-controlled transient elastography
(FibroScan) on a nationally representative sample in its 2017--March 2020 and
2021--2023 cycles \cite{nhanes_lux}. Both cycles use the same examination
protocol and variable names. Participants with a complete examination
(LUAXSTAT $=1$) and an observed controlled attenuation parameter (CAP) form
the cohort of 15{,}300 participants aged 12 to 80. A test set of 3{,}060
participants (20\%, stratified by cycle, seed 42) was set aside before any
model choice and is used only for evaluation; the remaining 12{,}240 form the
training set. The split is saved and reused by every later component, so the
same test participants are never seen in training anywhere in the system.

The liver-state model is intended for people aged 12 to 59, evaluated against
reliable examinations (CAP interquartile range below 40~dB/m); this is the
model's \emph{scope} (Chapter~\ref{chap:experiments}). The test set holds
1{,}236 in-scope participants, who are the held-out patients for every
evaluation in this thesis.

\subsection{Habitual diet from the 24-hour recall}

Each NHANES participant reports everything eaten on the previous day in a
24-hour dietary recall. The day-1 nutrient totals of both cycles are joined
with the USDA Food Patterns Equivalents (FPED), which give added sugars and
food-group amounts. USDA has released FPED totals only for 2017--March 2020;
for 2021--2023 they are derived with USDA's own method, from the individual
foods and the 2017--2020 food-level FPED values (about 99.6\% of the
recalled energy has a matching food code). From the recall the twin computes
the same composition measures it uses for recipes---energy, saturated fat and
sugars as a share of energy, and fibre per 1{,}000~kcal---and the Healthy
Eating Index 2020 \cite{shamswhite2023hei}.

A single day is a noisy measure of an individual's usual diet
\cite{dodd2006usual}, and people under-report what they eat
\cite{subar2003open}. It is only in the event of NHANES judging the recall as reliable, 
and also if the participant indicates that the day was normal, 
that the recall may be considered (500--3{,}500~kcal for females, 800--4{,}000~kcal for males). Even in this case, 
the recall can affect the ranking of diet by not more than one level 
(Section~\ref{sec:m_layer3}). New users who are not NHANES participants type
what they ate yesterday in plain words; each item is linked to a USDA SR
Legacy food, converted to grams and summed into the same totals.

\subsection{Recipe corpus}

Recipes come from RecipeDB \cite{batra2020recipedb}, which records each
recipe's ingredients with their quantities, instructions, preparation time,
servings, nutrition per serving, cuisine, diet labels and user ratings. A
sample of 20{,}000 recipes was collected through the RecipeDB interface. Only
main meals are candidates for recommendation: desserts, drinks, sauces,
dressings and spreads are excluded by category and title, as are dishes with
less than 200~kcal per serving. This leaves 10{,}108 candidate meals. Most
(71\%) are labelled ``Other'' rather than with a specific cuisine; the
largest labelled cuisines are European (1{,}152), Asian (637), Mexican (308),
Canadian (247) and American (228).

Every ingredient line is linked to a USDA SR Legacy food and its amount
converted to grams, which allows two measures that RecipeDB does not provide:
the degree of processing and the glycaemic load (Section~\ref{sec:m_layer4}).

\subsection{NHANES III: external cohort with long-term follow-up}

The external validation uses NHANES~III (1988--1994), in which an archived
gallbladder ultrasound was re-read for hepatic steatosis in 2009--2010
\cite{nhanes3_ultrasound}. Adults aged 20--59 with a gradable image form a
cohort of 10{,}682 participants, 20.7\% of whom had moderate or severe
steatosis (the recode recommended by the National Center for Health
Statistics). NHANES~III examination and laboratory variables were mapped onto
the model's inputs with unit conversions; 51 of the 81 inputs are available,
and the rest are treated as missing, which the model handles natively. The
public-use Linked Mortality File follows the participants to 31 December 2019
\cite{nchs2022mortality}: 10{,}671 are linked, with 2{,}173 deaths over a
median of 27.3 years. Its cause-of-death recode distinguishes heart disease,
stroke and cancer but has no liver-disease category.

% ======================================================================
\section{Layer 2: Liver-State Model and Clinical Safety Checks}
\label{sec:m_layer2}

\subsection{Liver-state model}

The liver twin's core is an XGBoost regressor \cite{chen2016xgboost} that
estimates CAP from 81 routine inputs---body measurements, standard blood
tests, blood pressure and a short questionnaire---and six derived scores.
Two classifiers estimate the probability that CAP is at least 248~dB/m
(``some liver fat'', steatosis grade S1 or above) and at least 280~dB/m
(``a lot of liver fat'', grade S3) \cite{karlas2017cap}; both are calibrated
with isotonic regression \cite{zadrozny2002calibration}. The experiments that
led to this model are reported in Chapter~\ref{chap:experiments}.

The calibrated probabilities, not the CAP estimate, are the twin's primary
state: with a held-out $R^2$ of 0.56 the individual CAP error (RMSE about
41~dB/m) is wider than the 32~dB/m between the S1 and S3 cut-offs, so a
single CAP number would suggest more precision than the model has. Each
probability is turned into a yes/no decision at the threshold that maximised
the F1 score on the training data (0.432 for CAP $\geq$ 248 and 0.365 for
CAP $\geq$ 280).

\subsection{Uncertainty and explanation}

Each patient's CAP estimate carries a 90\% conformal prediction interval
\cite{vovk2005conformal}: five-fold out-of-fold predictions on the training
set give one absolute error per in-scope training patient, and the interval
is the estimate plus or minus their 90th percentile, with the finite-sample
correction. Each calibrated probability carries a band, the 5th to 95th
percentile of the probabilities given by isotonic calibrators refitted on 200
bootstrap resamples. SHAP values \cite{lundberg2017shap}, computed with
XGBoost's built-in TreeSHAP, split each patient's CAP estimate into the
average prediction plus one contribution per input; they describe what drives
the model's estimate, not causes that diet can change.

\subsection{Fibrosis risk, liver label and safety notices}

Diet advice must not replace the clinical care pathway. Three rule-based
checks run alongside the model.

\textbf{Fibrosis risk.} FIB-4 $=$ age $\times$ AST $/$ (platelets $\times$
$\sqrt{\mathrm{ALT}}$) \cite{sterling2006fib4} is computed for adults and
categorised by the EASL--EASD--EASO 2024 pathway \cite{easl2024masld}: below
1.3 low risk, 1.3--2.67 indeterminate (further testing advised) and above
2.67 high risk. A measured FibroScan liver stiffness, when available and
reliable, replaces FIB-4 (below 8~kPa low, 8--12~kPa indeterminate, 12~kPa or
more high). Because FIB-4 is less accurate under 35 \cite{rinella2023aasld},
a low score in a young adult who is likely to have liver fat carries a note
that elastography would help. A high risk withholds every recipe until the
patient has had a clinical assessment.

\textbf{Liver label.} Following the 2023 multisociety nomenclature
\cite{rinella2023nomenclature}, steatotic liver disease with at least one
cardiometabolic risk factor is labelled MASLD; with an alcohol intake of
140--350~g a week (women) or 210--420~g (men) MetALD; above that range
alcohol-related liver disease; and without a risk factor cryptogenic.
Adolescents are checked against the nomenclature's paediatric criteria. This is simply a 
research label that is used in the explanations and not a diagnosis. 
\textbf{Safety notices.} Those below 18 years of age, patients who have consumed alcohol at or above the MetALD
 level, on glucose lowering drugs, or those who are pregnant will get a note asking to 
consult a doctor before making any changes to the diet (Section~\ref{sec:m_update}).

% ======================================================================
\section{Layer 3: Evidence Engine}
\label{sec:m_layer3}

\subsection{Intervention evidence table}

The evidence engine rests on a curated table of 25 intervention--outcome
findings. Each row records the study, intervention, outcome, reported effect
and direction, design, population, sample size, limitations and an explicit
evidence-strength label: 4 rows are STRONG (for example, among patients who
lost 5\% or more of their weight in a 12-month prospective study with paired
liver biopsies, 58\% had resolution of steatohepatitis
\cite{vilargomez2015weight}), 11 MODERATE, 2 GUIDELINE RECOMMENDATION, 5 WEAK
or INCONCLUSIVE and 3 OBSERVATIONAL ASSOCIATION. Observational associations
are recorded but never used as effects.

\subsection{From patient state to dietary priorities}

Layer~2's decisions and routine measurements set explicit yes/no flags with
published thresholds: probable steatosis and probable moderate-to-severe
steatosis (from the calibrated probabilities); overweight and obesity (BMI 25
and 30~kg/m$^2$, or 23 and 27.5 for Asian patients); a raised waist (above
102~cm in men and 88~cm in women); raised glycaemia (HbA1c 5.7\% or more); and
dyslipidaemia (raised triglycerides or low HDL). Evidence rules then fire on
these flags and produce up to six \emph{priorities}, each with a level (low,
medium or high) and the evidence rows that support it
(Table~\ref{tab:m_priorities}). Every priority requires probable steatosis;
the other flags decide which priorities apply and how high they rank.
Priorities say which dietary objective matters
more for this patient; they are not prescriptions and carry no numeric
intake target.

\begin{table}[htbp]
\centering
\caption{Dietary priorities: when each applies (always together with probable steatosis), its default level, what raises it, and the recipe criterion it sets.}
\label{tab:m_priorities}
\begin{tabular}{p{2.6cm}p{2.4cm}p{1.3cm}p{3.8cm}p{3.6cm}}
\toprule
Priority & Also requires & Default & Raised to & Recipe criterion (reference) \\
\midrule
Weight management & Excess weight & Medium & High: probable CAP $\geq$ 280, or obesity & Energy per serving $\leq$ 600~kcal \\
Diet quality & -- & Medium & High: probable CAP $\geq$ 280 & Fibre $\geq$ 14~g per 1{,}000~kcal \\
Limit saturated fat & -- & Medium & High: dyslipidaemia & Saturated fat $\leq$ 10\% of energy \\
Limit excess sugar & -- & Low & Medium: raised glycaemia or excess weight & Sugars $\leq$ 10\% of energy \\
Limit ultra-processed foods & -- & Low & -- & Processing FPro $\leq$ 0.5 \\
Lower glycaemic load & Raised glycaemia & Low & -- & Glycaemic load $\leq$ 20 per serving \\
\bottomrule
\end{tabular}
\end{table}

When the patient has a reliable recall, a priority moves up one level if the
patient's own intake misses the same reference value used for recipes, and
down one level if it meets it.

\subsection{Causal knowledge graph}

A small hand-written causal graph of 16 links encodes the well-established
direction of MASLD biology: dietary and activity levers act on body weight,
insulin resistance and liver fat; liver fat drives liver injury and fibrosis;
CAP, ALT and liver stiffness only measure them. Each link names its evidence
rows. The graph confines the prioritization to the levers having a causal pathway 
towards liver fat, and therefore the criteria of any recipe never focuses on any downstream 
measurement such as ALT, and it provides the chain for every recommendation for the 
explanations (Section~\ref{sec:m_layer4}).

% ======================================================================
\section{Layer 4: Recipe Evaluation and Ranking Engine}
\label{sec:m_layer4}

Layer~4, NutriTwin, is the core contribution of this thesis. It scores every
candidate meal on five criteria, ranks first the meals that meet every
evidence-linked reference value, and explains each recommendation.

\subsection{Food preferences}

At sign-up the patient states how they eat (anything, pescatarian,
vegetarian, vegetarian without eggs, or vegan), any of the nine major food
allergens named in US law plus gluten, any other allergy typed in, foods never
eaten (pork, beef or alcohol), disliked and liked ingredients, the longest
time they can spend cooking, and how many people they cook for. Diet,
allergies, foods avoided, dislikes and cooking time are hard rules: a recipe
that breaks one is never shown. A recipe must pass both RecipeDB's own diet
label and a check of its ingredient words and food groups, because the labels
contain errors (162 recipes labelled vegan list meat).

\subsection{Criteria}

\textbf{Health relevance.} Each priority maps to one recipe-composition
criterion with a reference value from public dietary guidance
(Table~\ref{tab:m_priorities}): energy per main meal \cite{phe2018calories},
saturated fat \cite{who2023satfat,dga2020}, sugars \cite{who2015sugars} and
fibre \cite{dga2020}. Two criteria are computed from the ingredients. The
degree of processing (FPro, 0 to 1) is the energy-weighted mean of the
ingredients' FPro, predicted by a re-implementation of FoodProX
\cite{menichetti2023processing}, a random forest that infers the NOVA
processing class of a food from its nutrient profile. The glycaemic load per
serving is the sum over ingredients of glycaemic index times available
carbohydrate divided by 100 \cite{atkinson2021gi}, with each carbohydrate food
taking the median glycaemic index of its matching entries in the 2021
international tables. Both are left blank when too few ingredients can be
linked. A criterion scores 100 at or past its reference value and falls
linearly to 0 over a tolerance; the health score is the mean of the patient's
criteria weighted by priority level.

\textbf{Cuisine compatibility.} A recipe labelled with the requested cuisine
scores 100, a declared related cuisine 75 and any other label 0. Because 71\%
of candidates are unlabelled, these are judged on their ingredients: the
typicality of a recipe for a cuisine is the mean share of that cuisine's
recipes that use each pair of the recipe's ingredients, scaled so that the
cuisine's own recipes at the 75th percentile score 100. Unlabelled recipes
score at most 50, so they fill gaps without displacing correctly labelled
ones.

\textbf{Practical feasibility.} The mean of three components, each 100 at the
minimum and 0 at a screening limit: total time (60 minutes), number of
ingredients (15) and number of instruction steps (12). The score is then
reduced by 25\% for each demanding technique named in the instructions
(deep-frying, kneading or proving dough, overnight chilling, marinating,
thermometer work, flamb\'eing) and by 10\% for each piece of special
equipment.

\textbf{Acceptability.} RecipeDB's user rating, shrunk towards the corpus
mean (4.91) with the weight of 20 ratings so that a recipe rated 5.0 by one
user does not outrank one rated 4.9 by hundreds, and scaled from 4.5 (0) to
5.0 (100).

\textbf{Personal fit.} The closeness of a recipe to what the patient already
eats and likes, with two parts averaged when both are known: the distance of
the recipe from the patient's 24-hour recall on the patient's own priorities
(meal energy, saturated fat, sugar and fibre, each in units of the
criterion's tolerance, capped at 1), and one minus the largest ingredient
overlap (Jaccard index) with a favourite dish or the liked ingredients.
The personal fit is equal to 100 times (one minus this distance); 
if there is no reliable recall or favorites then it is unknown. 
\subsection{Health-gated ranking}

The total score is the weighted average of the criteria which the recipe satisfies, 
where the weights are renormalized according to the criteria satisfied (Table~\ref{tab:m_weights}). 
The weighted scoring alone allows such a trade-off. The ranking is, in this case, 
\emph{health-gated}: recipes that fulfill all of the evidence-based reference values 
rank above the rest, and the weight score ranks recipes within these categories. Both 
processing and glycaemic load determine the score but are not 
health gated since the evidence behind them is less conclusive.

Inside the gated group every health score would be 100 and the patient's
priority levels would no longer matter. There, half of the health component
is instead the priority-weighted margin by which the recipe goes past the
reference values, so a patient whose saturated-fat priority is high is
offered recipes further below the saturated-fat reference. Personal fit is
applied only within the gated group, after ranking, so familiarity never
outranks health.

\begin{table}[htbp]
\centering
\caption{Criterion weights. Weights are renormalised over the criteria a recipe has; personal fit counts only for recipes that meet every reference value.}
\label{tab:m_weights}
\begin{tabular}{lrr}
\toprule
Criterion & Weight & Share when all are present \\
\midrule
Health relevance & 0.6 & 0.400 \\
Cuisine compatibility & 0.2 & 0.133 \\
Practical feasibility & 0.2 & 0.133 \\
Acceptability & 0.1 & 0.067 \\
Personal fit & 0.4 & 0.267 \\
\bottomrule
\end{tabular}
\end{table}

The personal-fit weight was chosen on the evaluation of
Section~\ref{sec:m_eval}: 0.4 made the top-ten lists of different patients
clearly different while every recommended recipe still met every reference
value (Chapter~ef{chap:results}). The other weights were set in the proposal and tested for
sensitivity over a grid of 36 combinations.

\subsection{Explanations, diversity and weekly plans}

Each recommended recipe is explained by following its criteria back through
the causal knowledge graph to liver fat, with the evidence rows and strength
of each link, and ending with the cuisine and practicality context that kept
the recipe appropriate. Every statement is built from the ranking engine's
own output.

A seven-day plan of two main meals a day is drawn greedily from the gated
recipes, practical ones first, after re-ranking by maximal marginal relevance
\cite{carbonell1998mmr} ($\lambda = 0.9$) so that near-duplicates do not fill
the list. The solution needs to meet the constraints on a weekly basis. This involves 
meeting the requirements of each meal and each pair of daily meals based on the reference 
values of the patient, the recipe appearing twice or not at all in the week and not consecutively, 
there being a difference in protein families in daily meals, 
and the absence of a dominant protein family in the week. If the cuisine
cannot fill the week, a favourite may appear twice, and then recipes from the
most similar cuisines (by ingredient profile) are borrowed and labelled as
such. A shopping list adds up one serving of each planned meal.

% ======================================================================
\section{Layer 5: Scenario Engine}
\label{sec:m_layer5}

The patient chooses between two scenarios. \emph{Scenario~A} takes the
recipes in the health-gated order, the healthiest available.
\emph{Scenario~B} takes recipes from the same safe pool ordered by closeness
to what the patient eats now, so the change is smaller. For each criterion
the engine compares the scenario's mean composition with the patient's own
recall and reports whether it improves, worsens or does not change
meaningfully, each direction carrying the evidence rows of its pathway:
calorie restriction (and, if sustained to 5\% weight loss, the STRONG
weight-loss evidence \cite{vilargomez2015weight}); saturated fat
\cite{luukkonen2018saturated}; sugar, labelled MODERATE only when the scenario
also lowers energy, because sugar harms the liver when it adds surplus energy
but not when exchanged isocalorically \cite{chiu2014fructose}; and fibre as a
marker of a Mediterranean-style pattern \cite{haigh2022mediterranean}. The
scenario's label is the strongest among the criteria that improve, and every
scenario is labelled as indirect evidence, since one recipe is not a dietary
pattern.

The scenario engine never asks the Layer~2 model to predict an intervention's
effect, because the model was trained on a cross-sectional survey, not on
interventions. Two intermediate steps between the scenarios are used only by
the updating loop, which can move a patient gradually from Scenario~B
towards Scenario~A as they settle in, or back when they struggle
(Section~\ref{sec:m_update}).

% ======================================================================
\section{Updating the Twin: Plan, Check-ins and Re-assessment}
\label{sec:m_update}

\subsection{The plan}

Patients cannot be expected to know how large a change to make, so the twin
suggests a plan at three paces, taken from the trial prescriptions: gentle
(half the Diabetes Prevention Program calorie cut for the patient's weight
band and 90 minutes of walking a week), standard (the DPP cut of
500--1{,}000~kcal a day and 150 minutes) \cite{dpp2002lifestyle} and ambitious
(the Look AHEAD calorie goal and 175 minutes) \cite{lookahead2006intervention}.
Walking adds (MET $-1$) $\times$ weight $\times$ hours of energy use at 3.8~MET.
Patients under 18 receive no calorie cut. A plan with a safety notice waits
for a clinician's approval, and only a clinician can set custom numbers.

\subsection{Predict, observe, compare, adjust}

At each monthly check-in the patient enters their weight and how many of the
planned meals they cooked and how much they walked. The twin then runs one
cycle of a predict--observe--compare--adjust loop.

\emph{Predict.} The forecast follows Hall's published rule
\cite{hall2011energy}: a sustained change in intake moves weight by about
1~kg per 100~kJ a day in the end, half of it in about a year. The plan's full
effect is scaled by a personal response factor $k = a \times r$, the product
of adherence $a$ (the share of the plan followed, from the patient's answers)
and the body's response $r$ (1 for the textbook response).

\emph{Compare.} The 90\% range of the forecast combines the uncertainty in
$k$ with ordinary weight drift, set from the spread of weight change in Look
AHEAD usual care and assumed to grow like a random walk. Progress is also
judged against the plan from the start: from month 1.5, a patient who has
lost less than half of the loss the plan typically gives is \emph{behind
plan}. It can be seen from the patient's responses that if the adherence level is low, 
then the adherence to the regimen has not taken place, and vice versa. 
\emph{Adjust.} Response variable $r$ is reassessed by Bayesian linear regression using the changes 
in weight between visits from a weakly informative prior ($r = 1.0 \pm 0.7$), with 
older values discounted due to adherence drift. The projection at 3, 6 and 12 months uses what has been learned, and
the liver projection re-runs Layer~2 with the projected weight, BMI and waist.
A plan changed by the clinician part-way through is handled as a dated
schedule of energy changes, so each month is judged against the plan in force
at that time.

\subsection{Tips and alerts}

Every message follows from one rule on the twin's output. The patient that did not adhere 
to the plan will be provided with more familiar recipes, while the one that adhered to 
the plan will be advised to continue following it. After two consecutive follow-ups of 
patients on the plan, the doctor will receive a notification of the points to consider: 
intake greater than stated, regular fluctuations in weight, and medical issues like newly 
prescribed medication or hyperthyroidism. The alert does not claim the
patient's body is slow: the in-silico evaluation showed that this cannot be
told apart from weight drift within a year (Chapter~ef{chap:results}). A patient on plan
with good adherence at two check-ins in a row is offered one step towards the
healthiest recipes. Clinicians also see alerts for plans awaiting approval,
indeterminate or high fibrosis risk, overdue check-ins and missing blood
tests.

% ======================================================================
\section{Plain-Language Summary}
\label{sec:m_summary}

The patient reads a short summary of the twin's output. It is produced solely 
based on the facts computed by the twin, with all numbers being rounded up as 
indicated: the probabilities of a few and a high amount of fat in the liver 
(never as 0\% or 100\%), a clear verdict based on the twin's binary decision, 
the liver classification, the outcome of fibrosis using published thresholds 
and the sources for those thresholds (with an indication where the score falls 
within 10\% of the threshold), any warnings about risks, and the importance of 
criteria in clear language with the number of studies supporting each, and the 
top recipes. Evidence codes, the CAP estimate and its interval are kept for the clinician's view.
Notices about changing diet appear only when a diet change is suggested.

A template writes the summary, and a language model can be plugged in
instead. Whatever writes it, an automatic check confirms that every number
and every evidence citation in the text appears in the twin's facts; if not,
the template summary is shown and the reason recorded.

% ======================================================================
\section{System and Interfaces}
\label{sec:m_system}

The twin's Layers~2--5 are wrapped in an HTTP backend (FastAPI) with an
SQLite database that keeps each patient's dated visits, plans, check-ins and
twin runs, so later visits can be compared with what the twin expected.
Patients sign in with a one-time code and see only their own record;
clinicians see the patients of their clinic. One twin run has two views. The
patient's view holds the summary, the chances of liver fat, the safety check,
the priorities, the recipes with their reasons and the two scenarios, but no
CAP estimate. The clinician's view holds the full output, including the CAP
estimate, its 90\% interval and the SHAP drivers.

The dashboard (Layer~7) presents this journey: a short sign-up, the patient's
results, the choice of pace and recipe style, a ``Today'' screen with one next
step and the day's meals, the weekly plan and shopping list, monthly
check-ins, and progress against the plan; and for the clinician, a list of
patients ordered by alerts and a patient view in which plans are approved or
changed. The liver visualisation (Layer~6) colours a three-dimensional liver
by the steatosis state from the calibrated probabilities, with a timeline
slider for today and three, six and twelve months, and labels every future
state as a simulated projection, not a clinical forecast.

% ======================================================================
\section{Evaluation Design}
\label{sec:m_eval}

Each component is evaluated against appropriate baselines on data it was not
built on. Unless stated, intervals are 95\% bootstrap percentile intervals
over participants, and results are unweighted.

\subsection{Liver-state model and safety checks}

The model is evaluated on the in-scope test participants by $R^2$ and RMSE
for CAP, and by ROC-AUC, Brier score and calibration for the two
classifiers, against the Fatty Liver Index \cite{bedogni2006fli} and the
Hepatic Steatosis Index \cite{lee2010hsi} on the same participants
(Chapter~\ref{chap:experiments}). The conformal intervals are checked by their
empirical coverage on the test set (target 90\%). FIB-4 is evaluated against
liver stiffness measured in the same FibroScan examination, with 8~kPa to rule
out and 12~kPa to rule in advanced fibrosis.

\subsection{Fairness audit}

The calibrated classifiers are scored separately by sex, race and ethnicity,
and age group (12--17, 18--39, 40--59) on the in-scope test participants:
ROC-AUC, Brier score, and the mean predicted probability against the observed
share (calibration in the large), with 1{,}000 bootstrap resamples within each
group. Groups with fewer than 30 participants are not reported. The share of
recommended recipes meeting each ranking criterion is compared across the same
groups.

\subsection{External and long-term validation on NHANES III}

On the NHANES~III cohort, the model's probabilities are compared with the
ultrasound grade (moderate or severe against normal or mild) by ROC-AUC,
alongside the Fatty Liver Index (on participants with GGT), the Hepatic
Steatosis Index, BMI and waist computed from the same record. This is a
different reference standard 30 years earlier, so it tests whether the model
ranks people correctly, not whether it reproduces CAP. For mortality,
participants are divided into quarters of the probability of CAP $\geq$ 280,
and death rates per 1{,}000 person-years are compared for all causes,
cardiovascular causes (heart disease or stroke) and cancer. Rate ratios
adjusted for age and sex come from Poisson regression with the logarithm of
person-years as offset, fitted by iteratively reweighted least squares, with
Wald intervals. The ultrasound grade and the FIB-4 category are analysed the
same way.

\subsection{Recipe ranking}

The ranking engine is evaluated on 300 randomly chosen in-scope test
participants with at least one recipe criterion. NHANES records no cuisine
preference, so each participant is given a cuisine at random among those with
at least 100 candidate recipes. The main measure is the share of each top-ten
list that meets all three criteria at once: every one of the patient's health
reference values, the patient's cuisine, and the practical limits (60 minutes,
15 ingredients, 12 steps). The proposed ranking is compared with the same
ranking without personal fit, without acceptability and with diversity
re-ranking; with weighted scoring without the health gate; with health-only
and nutrition-only rankings; and with random selection. Personalisation is
measured by the number of distinct top-ten lists and the mean overlap of the
top ten between patients of the same cuisine, with and without the 24-hour
recall. Weekly plans are checked against every weekly constraint, and the
sensitivity of the ranking to its weights is measured over 36 weight
combinations.

\subsection{Plain-language summaries}

For 300 in-scope test participants the twin runs end to end. The summaries
are scored on faithfulness (every number and citation present in the twin's
facts) and on readability (Flesch Reading Ease and Flesch--Kincaid grade
\cite{flesch1948readability}). The faithfulness check itself is tested by
planting errors in each summary as a language model might make them---a
number changed, a number added, a citation invented, which should all be
caught; a harmless rewording, which should pass; and a wrong word with no
number, which measures the check's blind spot. A random sample of ten
summaries was read by the project team, each beside the words its rules should
produce.

\subsection{In-silico evaluation of the updating twin}

NHANES examined each participant once, so the twin's updating and projections
cannot be tested on real follow-up. Following the design of in-silico trials
\cite{pappalardo2019insilico} and plasmode simulation
\cite{franklin2014plasmode}, real test patients provide month~0 and their next
twelve months are simulated; this addresses the verification part of digital
twin validation \cite{nasem2024digitaltwins}.

\textbf{Simulator.} Each of the 1{,}236 in-scope test patients follows usual
diet, Scenario~A ($-500$~kcal a day) and Scenario~B ($-250$~kcal a day) for
twelve months. A person-level adherence pattern (steady, slipping, regaining
or dropping out) scales the prescribed change. Monthly weight follows Hall's
model \cite{hall2011energy}, scaled by each person's body fat, with a random
weight drift independent of the diet. Waist, blood markers and blood pressure
change with weight loss by published effect sizes \cite{wing2011modest};
liver fat falls with weight loss and further while intake is reduced
\cite{lim2011reversal}; CAP follows liver fat, and liver enzymes follow the
change in liver fat \cite{koutoukidis2019weightloss,wong2013community}. Each
person's biology is drawn once, so the scenarios differ only in diet. The
simulator never uses the Layer~2 model, so the twin is not checked against its
own output.

\textbf{Observation layer.} The twin never sees the simulated truth. It sees
what a clinic would record: monthly weights and food logs, blood tests every
three months with within-person and laboratory variation from the EFLM
database \cite{eflm_bv}, CAP at months 0, 6 and 12 with measurement error,
food logs that under-report intake as in the OPEN study \cite{subar2003open},
and missed visits that become more likely when a patient's weight has risen.

\textbf{Calibration.} The settings that no trial reports directly (adherence
patterns and how they fade, response spread, weight drift) were tuned by
differential evolution so that virtual patients resembling each trial's
participants, drawn from the training set, reproduce published outcomes: Look
AHEAD at one year \cite{lookahead2006intervention}, the DPP at 24 weeks
\cite{dpp2002lifestyle} and Vilar-Gomez et al.\ at one year
\cite{vilargomez2015weight}. During each programme's intensive-contact phase
a slipping participant keeps full effort. Later outcomes (Look AHEAD at eight
years \cite{lookahead2014eightyear} and the DPP at about three years) were held
back to check the tuned simulator, and first-year regain in Look AHEAD
\cite{neiberg2012patterns} was used only to choose how many participants stop
the diet completely. HepatoTwin's scenarios use the adherence fitted to
Vilar-Gomez et al., routine advice to patients with fatty liver.

\textbf{Measures.} The twin runs month by month on the observed data with the
patient's answers derived from the true adherence (honest answers) and, as a
robustness check, inflated by each person's under-reporting share, or absent.
The performance of the model is judged by: predictions of the weight at months 3 
and 6, 3 and 6 months ahead, compared to the textbook pattern without learning 
and to the last carried forward weight (mean absolute error and coverage of the 
90\% range); how the prediction error varies with the learning process from month 
1 to month 9; the reasons given for the slow learning and warnings, compared to actual 
adherence and body reaction; how many months it takes to detect a patient falling
behind and liver fat falling below 280~dB/m; the probability of CAP $\geq$
280 at months 0, 6 and 12 against the true CAP (AUC, Brier score,
calibration), and its month-6 forecast of month 12; robustness to more missed
visits, scale noise, over-reported answers and missing answers; and whether
its projections order Scenario~A, Scenario~B and usual diet as the truth does.
These results verify that the twin can update, track and adapt; they do not
show that it predicts real patients, and almost all calibration trials were in
adults.